\section*{Part III: Cash-Flow Management and the Working-Capital Engine}

\subsection*{20. A Cash Forecast Begins with the Bank, Not with Revenue Projections}

The 1st row of a practical forecast is the \textbf{reconciled opening cash balance}. Subsequent rows list receipts and disbursements on dates they are likely to occur, not on dates they would be convenient. Start with committed payments: wages, rent, contractors, taxes, debt, infrastructure, contracted refunds, and unavoidable delivery costs. Then add customer collections according to the expected payment date and confidence level. Prospective sales belong in a separate, appropriately discounted layer.

Calculate daily or weekly ending cash: opening available cash plus actually anticipated receipts minus cash disbursements. Roll the forecast forward after each actual bank movement. A forecast that cannot reconcile with the bank is a storytelling device.

Use multiple horizons:
\begin{itemize}
\item \textbf{A 13-week rolling forecast} is a practical operating window for many small firms because it reveals payment timing, not because 13 is a mathematically optimal number.
\item \textbf{A 6-to-18-month scenario forecast} helps with staffing, financing, and strategic commitments.
\item \textbf{An annual operating plan} translates these into economics and resource decisions.
\end{itemize}
Each solves a different temporal problem.

\subsection*{21. Why the Forecast Should Be Direct Before It Is Sophisticated}

A useful early-stage weekly cash forecast has perhaps 15 lines: opening bank balance, high-confidence collections, uncertain collections, refunds, payroll, contractors, software, taxes, rent, loan service, planned capital expenditure, closing cash, minimum reserve, and resulting headroom. Add complexity only when it would change a decision.

The founder must resist the comfort of smooth averages. ``Customer payments are ₹4 lakh a month'' does not tell you whether ₹4 lakh arrives on the 2nd or the 29th. The obligations falling on the 10th may not wait. Likewise, annual salaries divided by 12 omit bonus months, severance, or payroll taxes if they apply.

A sound cash forecast has 3 visible properties: \textbf{timing} of each entry, \textbf{confidence} in receipt or payment, and \textbf{ownership} of the action required. ``Collect ₹2 lakh from College A in week 6'' is incomplete without the accepted invoice, the contractual due date, the college contact, and the person who will follow up. Finance is partly operational coordination.

\subsection*{22. A 13-Week Example: A Positive Forecast That Nearly Fails}

Consider a business with ₹5.5 lakh opening cash and the following \emph{aggregated} 13-week expectations. The periods are grouped for readability; actual management should use weekly rows.

\begin{table*}[htbp]
\centering
\small
\begin{tabularx}{\textwidth}{Xrrr}
\toprule
\textbf{Cash Movement (₹ lakh)} & \textbf{Weeks 1 to 4} & \textbf{Weeks 5 to 8} & \textbf{Weeks 9 to 13} \\
\midrule
Opening cash & 5.5 & 4.8 & 4.3 \\
Expected collections & 3.0 & 4.0 & 3.5 \\
Contracted and necessary payments & $-$3.7 & $-$4.5 & $-$4.7 \\
\midrule
\textbf{Ending cash (base case)} & \textbf{4.8} & \textbf{4.3} & \textbf{3.1} \\
\midrule
\emph{Delayed-collection scenario} & & & \\
If ₹3.5 lakh slips past week 13 & 4.8 & 4.3 & \textbf{$-$0.4} \\
\bottomrule
\end{tabularx}
\caption{13-week cash forecast showing how a single collection slip causes a cash deficit.}
\end{table*}

The company appears solvent at the end of the horizon in the base case (₹3.1 lakh). But suppose the ₹3.5 lakh receipt scheduled in the final period slips beyond week 13 while the payments remain due. The new ending cash becomes \textbf{negative ₹0.4 lakh} ($-$₹40,000). A solvent-looking forecast has become a near-term funding crisis from \emph{1 single collection delay}, even though the underlying contract remains profitable.

A founder should therefore keep a minimum-cash line and model the consequences of a delayed major receipt. The useful question is not, ``Does the base case look comfortable?'' It is, ``Which 2 plausible delays would create a shortage, and what can we do before that happens?''

\subsection*{23. Confidence-Weighted Forecasts Are Not Permission to Spend Probabilities}

Suppose the company expects a ₹4 lakh contract payment with a 70 percent subjective chance of arriving by Friday. Writing ₹2.8 lakh into the bank forecast as though the company could pay 70 percent of its salary bills with it is dangerous. The actual receipt will usually be nearer ₹4 lakh or zero, not ₹2.8 lakh.

Use scenarios or separate \textbf{committed}, \textbf{highly probable}, and \textbf{speculative} collections. In the conservative cash floor, include only funds whose timing and availability are sufficiently secure for the obligation in question. An expected-value forecast can inform planning; it should not conceal the discrete nature of payment failures.

The same caution applies to funding discussions. ``Investors are interested'' belongs in the strategic pipeline, not in the cash account. The founder's financial model should make the world less surprising, not make uncertainty disappear through fractional arithmetic.

\subsection*{24. Accounts Receivable: Revenue Waiting to Become Cash}

A receivable is a right to payment, not payment itself. It has at least 4 dimensions: \textbf{amount}, \textbf{due date}, \textbf{probability of collection}, and \textbf{administrative friction}. 1 invoice may be unchallenged and 5 days from settlement. Another may depend on an unapproved change order. Their nominal values should not receive the same liquidity treatment.

Maintain an \textbf{aged receivables register}:
\begin{itemize}
\item Current (not yet due),
\item 1 to 30 days overdue,
\item 31 to 60 days overdue,
\item 61 to 90 days overdue, and
\item More than 90 days overdue.
\end{itemize}
Record the person responsible, disputed items, contractual evidence, collection notes, and next action. Do not assume receivables will be collected just because they are recognized in accounting.

A founder may be reluctant to follow up on institutional invoices for fear of disturbing the relationship. The better strategy is to design payment clarity \textbf{before delivery}: purchase order, approved contracting entity, invoice recipient, tax particulars, acceptance evidence, payment terms, and escalation route. Professional collection is a service process, not a personal confrontation.

\subsection*{25. DSO: How Slow Collection Changes Financing Requirements}

\textbf{Days sales outstanding (DSO)} is commonly estimated as average trade receivables divided by credit sales, multiplied by the number of days in the period:
\begin{equation}
\text{DSO} = \frac{\text{Average Receivables}}{\text{Total Credit Sales}} \times \text{Days in Period}.
\end{equation}
The result is an approximate collection-time indicator; it is sensitive to seasonality, sales growth, and the definition of receivables and credit sales.

Suppose a B2B training provider makes ₹24 lakh of credit sales over a 90-day quarter and average eligible receivables are ₹8 lakh:
\begin{equation}
\text{DSO} = \frac{₹8,00,000}{₹24,00,000} \times 90 = \mathbf{30\text{ days}}.
\end{equation}
If collection slows to 60 days at similar sales, average receivables roughly double to ₹16 lakh, tying up approximately another ₹8 lakh of working capital.

This is why a seemingly minor payment-term change can dominate a high-margin contract. A company with a tiny reserve may reasonably prefer a ₹12 lakh contract with sensible milestones over a ₹16 lakh contract paid entirely after completion. The choice is not simply margin versus margin; it is \textbf{cash exposure and probability of settlement}.

\subsection*{26. Inventory and the Cash Conversion Cycle}

For physical businesses, inventory can absorb vast amounts of cash before any sale. The standard operating heuristic is:
\begin{equation}
\text{Cash Conversion Cycle (CCC)} = \text{DIO} + \text{DSO} - \text{DPO},
\end{equation}
where DIO is days inventory outstanding and DPO is days payables outstanding. It approximates how long cash invested in inventory and receivables remains tied up, net of supplier credit.

Suppose an appliance distributor holds 45 days of inventory, collects from retailers after 35 days, and pays manufacturers after 30 days. Its approximate cash conversion cycle is:
\begin{equation}
\text{CCC} = 45 + 35 - 30 = \mathbf{50\text{ days}}.
\end{equation}
Faster growth in such a model requires significant financing even when every unit has positive gross margin. Reduce inventory days, accelerate collections, or negotiate appropriate supplier terms, and the cash burden changes.

Do not apply the formula mechanically to every service business. A prepaid online education company has little physical inventory and may receive customer cash before delivery. The economic advantage of that arrangement is real; the corresponding \textbf{deferred delivery obligation} is also real. The absence of inventory does not mean the absence of working-capital risk.
